%% chapters/02-tinjauan-pustaka.tex
%% Bab II — Tinjauan Pustaka
%% Catatan: Tinjauan Pustaka != Dasar Teori. Bab ini menyajikan EVOLUSI
%% pekerjaan terdahulu secara kronologis dan menempatkan disertasi ini
%% terhadap state-of-the-art. Aturan: lihat .cursor/rules/07-bab2-tinjauan-pustaka.mdc

\chapter{Tinjauan Pustaka}
\label{bab:tinjauan-pustaka}

Bab ini menyajikan tinjauan kritis atas pekerjaan terdahulu yang
relevan dengan disertasi ini. Pembahasan disusun secara kronologis
dari metode prediksi RUL klasik berbasis fisika hingga pendekatan
\emph{deep learning} mutakhir, serta menempatkan kontribusi disertasi
ini terhadap \emph{state-of-the-art}.

\section{Domain Prognostik Bantalan Gelinding}
\label{subsec:bab2_domain}

% Konteks sejarah singkat: ISO standards, condition monitoring, evolusi
% dari corrective ke predictive maintenance. Frekuensi karakteristik
% (BPFO/BPFI/BSF/FTF) sebagai fondasi diagnostik. Fitur statistik
% klasik (RMS, kurtosis) dalam pemantauan kondisi.
\dots

\section{Evolusi Metode Prediksi RUL}
\label{subsec:bab2_evolusi}

\subsection{Metode Berbasis Model Fisika}
\label{subsec:bab2_fisika}

% Paris law, Forman law, model akumulasi kerusakan. Kelebihan: dapat
% diinterpretasi. Keterbatasan: butuh model fisika spesifik bantalan
% dan kondisi operasi yang akurat.
\dots

\subsection{Metode Berbasis Data — Pra Era \emph{Deep Learning}}
\label{subsec:bab2_data_pre_dl}

% Hidden Markov Model, Particle Filter, SVR, Kalman filter, ARMA, dan
% kombinasi feature engineering manual + classifier dangkal.
\dots

\subsection{Era \emph{Deep Learning}: RNN, LSTM, dan GRU}
\label{subsec:bab2_dl_lstm}

% Pekerjaan-pekerjaan kunci 2016--2020 yang menerapkan LSTM/GRU pada
% PHM2012 dan dataset bearing lainnya. Pencapaian dan keterbatasan.
\dots

\subsection{Era Transformer dan Variannya}
\label{subsec:bab2_transformer}

% Self-attention untuk time series, Informer, Autoformer, FEDformer,
% PatchTST. Trade-off akurasi vs biaya komputasi.
\dots

\section{State-of-the-Art Terkini: xLSTM dan Mamba}
\label{subsec:bab2_sota}

\subsection{xLSTM dan \emph{Matrix Memory}}
\label{subsec:bab2_xlstm}

% Pekerjaan Beck dkk. (2024) yang memperkenalkan xLSTM dengan komponen
% sLSTM dan mLSTM. Aplikasi xLSTM pada time series.
\dots

\subsection{Mamba dan \emph{Selective State-Space Models}}
\label{subsec:bab2_mamba}

% Pekerjaan Gu \& Dao (2023) tentang Mamba. Adopsi pada time series
% forecasting dan prognostik.
\dots

\section{Interpretabilitas pada Model \emph{Deep Learning}}
\label{subsec:bab2_interpret}

% SHAP, LIME, Integrated Gradients pada model time series. Sparse
% Autoencoder di NLP/vision (Anthropic features, OpenAI scaling).
% Adopsi terbatas SAE pada prognostik.
\dots

\section{Posisi Disertasi dan Celah Penelitian}
\label{subsec:bab2_posisi}

Berdasarkan tinjauan di atas, dapat diidentifikasi tiga celah utama
yang akan dialamatkan oleh disertasi ini:

\begin{enumerate}
  \item Belum ada arsitektur yang secara eksplisit menggabungkan
        \emph{state-space model} selektif dengan \emph{matrix-memory
        recurrent} untuk prediksi RUL bantalan.
  \item Aplikasi \emph{Sparse Autoencoder} pada prognostik bantalan
        masih sangat terbatas, dan belum ada studi yang memetakan
        fitur SAE ke frekuensi karakteristik fisik.
  \item Studi generalisasi lintas dataset masih jarang, sehingga
        klaim SOTA pada satu dataset belum tentu konsisten lintas
        kondisi operasi.
\end{enumerate}

% Tabel komparatif: 8-12 baris pekerjaan terkait, kolom: tahun, metode,
% dataset, metrik, kontribusi, keterbatasan. Letakkan baris terakhir
% sebagai disertasi ini.
\begin{table}[ht]
\centering
\caption{Posisi disertasi terhadap pekerjaan terdahulu yang representatif.}
\label{tab:bab2_posisi}
\small
\begin{tabular}{@{}p{2.2cm}p{2cm}p{2cm}p{2.2cm}p{4cm}@{}}
\toprule
Studi & Arsitektur & Dataset & Interpret. & Keterbatasan \\
\midrule
\dots & \dots & \dots & \dots & \dots \\
\midrule
\textbf{Disertasi ini} & Mamba-xLSTM-Net & PHM2012, XJTU-SY & SAE + SHAP + IG & --- \\
\bottomrule
\end{tabular}
\end{table}
